\textbf{Program 1: Descending quick sort with trace and counts.}

\textbf{File:} \texttt{quick\_sort\_descending.cpp}.

\textbf{Chapter connection:} Quick sort.

Move the middle-index pivot to the end, partition larger values to its left, then recurse on the two remaining ranges. The trace highlights active indexes. Counted swaps and key comparisons are reported consistently across the test cases.

\begin{adslabfigure}[How this program runs]
\begin{adslabgraphic}
\begin{tikzpicture}[font=\small,execute at begin node={\hyphenpenalty=10000\exhyphenpenalty=10000}]
\node[utpstep,text width=2.50cm,minimum height=1.15cm] (s0) at (0.0,0) {\textbf{1}\\Choose middle pivot};
\node[utpstep,text width=2.50cm,minimum height=1.15cm] (s1) at (2.95,0) {\textbf{2}\\Partition descending};
\node[utpstep,text width=2.50cm,minimum height=1.15cm] (s2) at (5.9,0) {\textbf{3}\\Recurse left / right};
\node[utpstep,text width=2.50cm,minimum height=1.15cm] (s3) at (8.85,0) {\textbf{4}\\Verify and count};
\draw[utparrow] (s0.east)--(s1.west);
\draw[utparrow] (s1.east)--(s2.west);
\draw[utparrow] (s2.east)--(s3.west);
\end{tikzpicture}
\end{adslabgraphic}
\end{adslabfigure}


\begin{utpsteps}\setcounter{enumi}{0}\item \textbf{Headers and initial definitions.}\end{utpsteps}

\texttt{quick\_sort\_descending.cpp} --- lines 17--26.

\lstinputlisting[firstline=17,lastline=26]{all_experiments/Lab10/quick_sort_descending.cpp}


\tcbbreak

\begin{utpsteps}\setcounter{enumi}{1}\item \textbf{Implement show, countedSwap, partitionDescending.}\end{utpsteps}

\texttt{quick\_sort\_descending.cpp} --- lines 27--61.

% Worked lab input; chapter fonts, colours and TikZ styles retained.
\begin{adslabfigure}[Descending partition: trace the entire first call]
\begin{adslabgraphic}
\begin{tikzpicture}[font=\small]
\node[utpcell,minimum width=0.88cm,minimum height=.65cm] at (0.00,0.00) {29};
\node[utplabel] at (0.00,0.60) {0};
\node[utpcell,minimum width=0.88cm,minimum height=.65cm] at (0.92,0.00) {10};
\node[utplabel] at (0.92,0.60) {1};
\node[utphot,minimum width=0.88cm,minimum height=.65cm] at (1.84,0.00) {14};
\node[utplabel] at (1.84,0.60) {2};
\node[utpcell,minimum width=0.88cm,minimum height=.65cm] at (2.76,0.00) {37};
\node[utplabel] at (2.76,0.60) {3};
\node[utpcell,minimum width=0.88cm,minimum height=.65cm] at (3.68,0.00) {13};
\node[utplabel] at (3.68,0.60) {4};
\node[utplabel,anchor=west,align=left,text width=5.2cm] at (4.8,0.00) {Middle index 2: pivot is 14};
\draw[utparrow] (-.85,-0.35)--(-.85,-0.70);
\node[utpcell,minimum width=0.88cm,minimum height=.65cm] at (0.00,-1.05) {29};
\node[utpcell,minimum width=0.88cm,minimum height=.65cm] at (0.92,-1.05) {10};
\node[utpcell,minimum width=0.88cm,minimum height=.65cm] at (1.84,-1.05) {13};
\node[utpcell,minimum width=0.88cm,minimum height=.65cm] at (2.76,-1.05) {37};
\node[utphot,minimum width=0.88cm,minimum height=.65cm] at (3.68,-1.05) {14};
\node[utplabel,anchor=west,align=left,text width=5.2cm] at (4.8,-1.05) {\texttt{swap(a[mid],a[high])}; set i=-1};
\draw[utparrow] (-.85,-1.40)--(-.85,-1.75);
\node[utphot,minimum width=0.88cm,minimum height=.65cm] at (0.00,-2.10) {29};
\node[utpcell,minimum width=0.88cm,minimum height=.65cm] at (0.92,-2.10) {10};
\node[utpcell,minimum width=0.88cm,minimum height=.65cm] at (1.84,-2.10) {13};
\node[utpcell,minimum width=0.88cm,minimum height=.65cm] at (2.76,-2.10) {37};
\node[utpcell,minimum width=0.88cm,minimum height=.65cm] at (3.68,-2.10) {14};
\node[utplabel,anchor=west,align=left,text width=5.2cm] at (4.8,-2.10) {j=0: 29 $\ge$ 14; increment i to 0};
\draw[utparrow] (-.85,-2.45)--(-.85,-2.80);
\node[utpcell,minimum width=0.88cm,minimum height=.65cm] at (0.00,-3.15) {29};
\node[utphot,minimum width=0.88cm,minimum height=.65cm] at (0.92,-3.15) {10};
\node[utphot,minimum width=0.88cm,minimum height=.65cm] at (1.84,-3.15) {13};
\node[utpcell,minimum width=0.88cm,minimum height=.65cm] at (2.76,-3.15) {37};
\node[utpcell,minimum width=0.88cm,minimum height=.65cm] at (3.68,-3.15) {14};
\node[utplabel,anchor=west,align=left,text width=5.2cm] at (4.8,-3.15) {j=1,2: 10 and 13 are smaller; skip both};
\draw[utparrow] (-.85,-3.50)--(-.85,-3.85);
\node[utpcell,minimum width=0.88cm,minimum height=.65cm] at (0.00,-4.20) {29};
\node[utphot,minimum width=0.88cm,minimum height=.65cm] at (0.92,-4.20) {37};
\node[utpcell,minimum width=0.88cm,minimum height=.65cm] at (1.84,-4.20) {13};
\node[utphot,minimum width=0.88cm,minimum height=.65cm] at (2.76,-4.20) {10};
\node[utpcell,minimum width=0.88cm,minimum height=.65cm] at (3.68,-4.20) {14};
\node[utplabel,anchor=west,align=left,text width=5.2cm] at (4.8,-4.20) {j=3: 37 $\ge$ 14; increment i to 1\\Swap indexes 1 and 3};
\draw[utparrow] (-.85,-4.55)--(-.85,-4.90);
\node[utpcell,minimum width=0.88cm,minimum height=.65cm] at (0.00,-5.25) {29};
\node[utpcell,minimum width=0.88cm,minimum height=.65cm] at (0.92,-5.25) {37};
\node[utphot,minimum width=0.88cm,minimum height=.65cm] at (1.84,-5.25) {14};
\node[utpcell,minimum width=0.88cm,minimum height=.65cm] at (2.76,-5.25) {10};
\node[utpcell,minimum width=0.88cm,minimum height=.65cm] at (3.68,-5.25) {13};
\node[utplabel,anchor=west,align=left,text width=5.2cm] at (4.8,-5.25) {Swap pivot with a[i+1]; return p=2\\Recurse on indexes 0--1 and 3--4};
\end{tikzpicture}
\end{adslabgraphic}
\end{adslabfigure}


\lstinputlisting[firstline=27,lastline=61]{all_experiments/Lab10/quick_sort_descending.cpp}


\tcbbreak

\begin{utpsteps}\setcounter{enumi}{2}\item \textbf{Connect the test cases in main.}\end{utpsteps}

\texttt{quick\_sort\_descending.cpp} --- lines 62--104.

\lstinputlisting[firstline=62,lastline=104]{all_experiments/Lab10/quick_sort_descending.cpp}


\tcbbreak

\textbf{Build, run, and inspect quick\_sort\_descending.}

From this program\textquotesingle s folder, run these commands in PowerShell:

\begin{lstlisting}[style=utpterminal,language={}]

g++ -std=c++17 -Wall -Wextra quick_sort_descending.cpp -o quick_sort_descending.exe

.\quick_sort_descending.exe

\end{lstlisting}

\textbf{Expected output excerpt.} Compare values, order, and pass/fail messages; spacing need not match.

\begin{lstlisting}[style=utpterminal,language={}]

== Trace on the common dataset ==
  start: [29 10 14 37 13]   ([ ] marks the range being partitioned)
  pivot 14 (middle of 0..4) -> [29 37 14 10 13]   pivot now at index 2
  pivot 29 (middle of 0..1) -> [37 29] 14 10 13   pivot now at index 1
  pivot 10 (middle of 3..4) -> 37 29 14 [13 10]   pivot now at index 4
  result: [37 29 14 13 10]
  comparisons: 6, swaps: 5
  [ok]   the result is in descending order

== More inputs ==
  [ok]   duplicates: 5 3 5 1 5 9 3
  [ok]   already ascending: 1 2 3 4 5 6
  [ok]   already descending: 6 5 4 3 2 1
  [ok]   all equal: 7 7 7 7
  [ok]   one element and empty input

\end{lstlisting}

The complete captured run is in \texttt{labs/outputs/quick\_sort\_descending.txt}.

\begin{tryit}

Use empty, one-value, all-equal, and reversed arrays. Predict which ranges end immediately, and verify descending order without losing duplicate values.

\end{tryit}
